%% kapa: exam=1 round=37 year=2026 period=2 form=A subject=accounting subject_ko=회계학 seq=2 number=42
%% source: Q-Net 감정평가사 기출문제 (cst00309) artlSeq=5258924 file=raw/1008/5258924/2026년 제37회 감정평가사 1차 2교시.pdf
%% conversion: kordoc 4.24.51 -> tools/gen1.py (auto); review=proofread-sonnet; flags=none
\begin{question}{42}
재무제표 표시에 관한 설명으로 옳지 않은 것은?
\begin{choices}
  \item 부적절한 회계정책은 이에 대하여 공시나 주석 또는 보충자료를 통해 설명하더라도 정당화될 수 없다.
  \item 보고기간말 이전에 장기차입약정의 약정사항을 위반했을 때 대여자가 즉시 상환을 요구할 수 있는 채무는 보고기간 후 재무제표 발행승인일 전에 대여자가 약정위반을 이유로 상환을 요구하지 않기로 합의한 경우 비유동부채로 분류한다.
  \item 기업은 현금흐름 정보를 제외하고는 발생기준 회계를 사용하여 재무제표를 작성한다.
  \item 재무제표가 계속기업의 기준하에 작성되지 않는 경우에는 그 사실과 함께 재무제표가 작성된 기준 및 그 기업을 계속기업으로 보지않는 이유를 공시하여야 한다.
  \item 한국채택국제회계기준에서 요구하거나 허용하지 않는 한 자산과 부채 그리고 수익과 비용은 상계하지 아니한다.
\end{choices}
\end{question}
\answer{2}
